% from ega3-1-fr.tex
\oldpage[III]{157}

(4.6.8) en notant comme ci-dessus que $S$ est le seul voisinage de $S'$ dans $S$.
Soit $g:X\to T$ l'isomorphisme réciproque de $p\circ i$, et posons
$f=q\circ i\circ g$, qui est un morphisme $X\to Y$ dont par définition le graphe
$\Gamma_f=i\circ g$. Comme $\widehat g$ est l'isomorphisme réciproque de
$\widehat{p\circ i}=w$, on a
\[
\widehat{\Gamma_f}=\widehat i\circ\widehat g=j\circ w=\Gamma_h.
\]
Mais on sait que $\widehat{\Gamma_f}=\Gamma_{\widehat f}$ (I, 10.9.8), d'où
finalement $h=\widehat f$, ce qui achève la démonstration.
\end{proof}

On peut donc dire, dans le langage des catégories, que le foncteur
$X\mapsto\widehat X$ est pleinement fidèle ($0$, 8.1.6) de la catégorie des
schémas propres sur $\operatorname{Spec}(A)$ dans la catégorie des schémas formels
propres sur $\operatorname{Spf}(A)$, pour tout anneau adique noethérien $A$~; il
établit par suite une équivalence entre la première de ces catégories et une
sous-catégorie de la seconde~; les objets de cette dernière seront appelés
\emph{schémas formels algébrisables}. Pour un tel schéma $\mathfrak X$, il existe
un schéma usuel $X$, propre sur $\operatorname{Spec}(A)$, déterminé à isomorphisme
unique près, tel que $\mathfrak X$ soit isomorphe à $\widehat X$.

\begin{env}[5.4.2]
\label{III.5.4.2-fr}
\emph{Scholie.}
Avec les notations de (\hyperref[III.5.4.1-fr]{5.4.1}), posons
$S_n=\operatorname{Spec}(A/\mathfrak J^{n+1})$,
$X_n=X\times_S S_n$, $Y_n=Y\times_S S_n$. Il résulte de
(\hyperref[III.5.4.1-fr]{5.4.1}) et de (I, 10.12.3) que se donner un
$S$-morphisme $f:X\to Y$ équivaut à se donner un $(S_n)$-système inductif adique
(I, 10.12.2) de $S_n$-morphismes $f_n:X_n\to Y_n$.
\end{env}

\begin{remark}[5.4.3]
\label{III.5.4.3-fr}
Contrairement à ce que pourrait suggérer le th. d'existence
(\hyperref[III.5.1.6-fr]{5.1.6}), il y a des schémas formels propres sur
$\operatorname{Spf}(A)$ et qui ne sont pas algébrisables (tout comme il y a des
espaces analytiques compacts qui ne proviennent pas de variétés algébriques
complexes). Nous rencontrerons plus tard de tels schémas dans la «~théorie des
modules~», qui traite précisément (lorsque le corps de base est $\mathbb C$) des
variations infinitésimales de la structure complexe d'une variété algébrique
complète, et on sait que de telles variations peuvent donner naissance à des
variétés analytiques qui ne sont pas algébriques.
\end{remark}

\begin{proposition}[5.4.4]
\label{III.5.4.4-fr}
Soient $A$ un anneau adique noethérien, $\mathfrak S=\operatorname{Spf}(A)$,
$g:\mathfrak X\to\mathfrak S$, $h:\mathfrak Y\to\mathfrak S$ deux morphismes
propres de schémas formels, $f:\mathfrak X\to\mathfrak Y$ un
$\mathfrak S$-morphisme. Si $f$ est fini et si $\mathfrak Y$ est algébrisable,
alors $\mathfrak X$ est algébrisable.
\end{proposition}

\begin{proof}
On notera que les hypothèses sur $g$ et $h$ entraînent déjà que $f$ est propre
(3.4.1), et pour que $f$ soit fini, il suffit que pour tout $y\in\mathfrak Y$, la
fibre $f^{-1}(y)$ soit finie (\hyperref[III.4.8.11-fr]{4.8.11}). L'hypothèse
entraîne que $\mathcal B=f_*(\mathcal O_{\mathfrak X})$ est une
$\mathcal O_{\mathfrak Y}$-Algèbre cohérente
(\hyperref[III.4.8.6-fr]{4.8.6}), donc il résulte du th. d'existence que, si
$\mathfrak Y=\widehat Y$ et $h=\widehat w$, où
$w:Y\to\operatorname{Spec}(A)$ est un morphisme propre de schémas usuels, il
existe une $\mathcal O_Y$-Algèbre cohérente $\mathcal C$ telle que
$\mathcal B=\widehat{\mathcal C}$. Soient $X=\operatorname{Spec}(\mathcal C)$, et
$u:X\to Y$ le morphisme structural~; alors, il résulte aussitôt de la définition de
$\mathfrak X$ à partir de $\mathcal B$
(\hyperref[III.4.8.7-fr]{4.8.7}) que $\mathfrak X$ est canoniquement isomorphe à
$\widehat X$ et que $f$ s'identifie à $\widehat u$ (il suffit pour le voir de
considérer le cas où $Y$ est affine).

On notera que (\hyperref[III.5.1.8-fr]{5.1.8}) est un cas particulier de
(\hyperref[III.5.4.4-fr]{5.4.4}).
\end{proof}

\begin{theorem}[5.4.5]
\label{III.5.4.5-fr}
Soient $A$ un anneau adique noethérien, $\mathfrak J$ un idéal de définition de
$A$, $S=\operatorname{Spec}(A)$,
$\mathfrak S=\widehat S=\operatorname{Spf}(A)$,
$f:\mathfrak X\to\mathfrak S$ un morphisme propre de schémas formels. On pose
$S_k=\operatorname{Spec}(A/\mathfrak J^{k+1})$,
$X_k=\mathfrak X\times_{\mathfrak S}S_k$, et pour tout
$\mathcal O_{\mathfrak X}$-Module $\mathcal F$,
$\mathcal F_k=\mathcal F\otimes_{\mathcal O_{\mathfrak X}}\mathcal O_{X_k}
=\mathcal F/\mathfrak J^{k+1}\mathcal F$.

